\documentclass[11pt,a4paper]{article}

% ---------------------------------------------------------
% PACKAGES
% ---------------------------------------------------------
\usepackage[utf8]{inputenc}
\usepackage[T1]{fontenc}
\usepackage{lmodern}
\usepackage{amsmath, amssymb, amsfonts, bm}
\usepackage{geometry}
\usepackage{graphicx}
\usepackage{hyperref}
\usepackage{authblk}
\usepackage{cite}
\usepackage{physics}
\usepackage{mathtools}
\usepackage{tensor}
\usepackage{enumitem}

\geometry{margin=2.4cm}

% ---------------------------------------------------------
% TITLE
% ---------------------------------------------------------
\title{\textbf{Companion IV (v49.1)}\\[0.2cm]
\textbf{Thermodynamic Sector and Dissipation Blockade}\\[0.2cm]
\large Foundations of Relaxation-Driven Cyclic Cosmology (RDCC)}

\author[1]{Michael Lehmann}
\affil[1]{\small Independent Researcher, Relaxation-Driven Cyclic Cosmology (RDCC)\\
\texttt{mi.lehmann@gmx.de}}

\date{September 2026}

\begin{document}
\maketitle

% ========================================================
% ABSTRACT
% ========================================================

\begin{abstract}
This Companion develops the thermodynamic sector of RDCC in its fully
reconstructed and Phase-1-integrated v49.1 form. The visible and mirror sectors
possess opposite thermodynamic arrows of time, ensuring global CPT invariance.
The mirror sector cannot dissipate energy, leading to a thermodynamic blockade
that prevents the formation of compact objects and results in diffuse,
non-compact halos. We derive the global entropy structure, the Lindblad-like
effective dissipator, and the sectoral projection of thermodynamic quantities.
These results integrate the dimension-8 inter-sector operator, the global
definition of $\alpha(t)$, and the ultraviolet structure developed in
Companions III and X.
\end{abstract}

\tableofcontents
\newpage
% ========================================================
% 1 INTRODUCTION
% ========================================================

\section{Introduction}

Thermodynamics plays a central role in RDCC. The Universe is described by a
global CPT-symmetric quantum state defined on a bipartite Hilbert space,
\begin{equation}
\mathcal{H} = \mathcal{H}_{(+)} \otimes \mathcal{H}_{(-)}.
\end{equation}

The visible and mirror sectors arise as reduced density matrices of this global
state and possess thermodynamic arrows of time that point in opposite
directions. This inversion ensures global CPT invariance even though each
sector individually exhibits irreversible behaviour.

In standard cosmology, dissipation is possible in all sectors, allowing the
formation of compact objects such as stars, galaxies, and black holes. In RDCC,
dissipation is blocked in the mirror sector due to its reversed thermodynamic
arrow. As a consequence, the mirror sector cannot form compact objects and
instead forms diffuse, non-compact halos.

This Companion develops the corrected thermodynamic sector of RDCC in its
v49.1 form. We derive the global entropy structure, establish the thermodynamic
dissipation blockade, and present the Lindblad-like effective dissipator that
governs the irreversible behaviour of the visible sector.
% ========================================================
% 2 THERMODYNAMIC ARROWS AND GLOBAL ENTROPY
% ========================================================

\section[Thermodynamic Arrows and Global Entropy]{Thermodynamic Arrows and Global Entropy}

The global state is written as
\begin{equation}
\ket{\Psi} = \sum_n \lambda_n \ket{n}_{(+)} \otimes \ket{n}_{(-)},
\end{equation}
with reduced density matrices
\begin{align}
\rho_{(+)} &= \Tr_{(-)} \ket{\Psi}\bra{\Psi}, \\
\rho_{(-)} &= \Tr_{(+)} \ket{\Psi}\bra{\Psi}.
\end{align}

\subsection[Thermodynamic Arrow of the Visible Sector]{Thermodynamic Arrow of the Visible Sector}

The visible sector evolves forward in time,
\begin{equation}
\frac{d}{dt} S_{(+)}(t) \ge 0,
\end{equation}
where $S_{(+)} = -\Tr(\rho_{(+)} \log \rho_{(+)})$ is the von Neumann entropy.

\subsection[Thermodynamic Arrow of the Mirror Sector]{Thermodynamic Arrow of the Mirror Sector}

The mirror sector evolves backward in time,
\begin{equation}
\frac{d}{dt} S_{(-)}(-t) \ge 0.
\end{equation}

This inversion ensures global CPT invariance,
\begin{equation}
S_{(+)}(t) + S_{(-)}(-t) = \text{constant}.
\end{equation}

\subsection[Global Entropy Conservation]{Global Entropy Conservation}

The global entropy
\begin{equation}
S_{\mathrm{global}} = -\Tr(\ket{\Psi}\bra{\Psi} \log \ket{\Psi}\bra{\Psi})
\end{equation}
is constant because the global state is pure. Irreversibility arises only after
sectoral projection.
% ========================================================
% 3 DISSIPATION BLOCKADE AND LINDBLAD STRUCTURE
% ========================================================

\section[Dissipation Blockade and Lindblad Structure]{Dissipation Blockade and Lindblad Structure}

The mirror sector cannot dissipate energy. Its reversed thermodynamic arrow
prevents the formation of compact objects. Dissipation requires
\begin{equation}
\frac{d}{dt} S \ge 0,
\end{equation}
but the mirror sector evolves according to
\begin{equation}
\frac{d}{dt} S_{(-)}(-t) \ge 0,
\end{equation}
which forbids standard cooling processes.

\subsection[Non-Compact Halos]{Non-Compact Halos}

The mirror sector forms diffuse, non-compact halos. These halos contribute to
the gravitational potential experienced by the visible sector and are
interpreted as dark matter.
\subsection[Lindblad-Like Dissipator]{Lindblad-Like Dissipator}

The visible sector evolves according to
\begin{equation}
\dot{\rho}_{(+)} =
- i [H_{(+)}, \rho_{(+)}]
+ \mathcal{D}_{\mathrm{eff}}[\rho_{(+)}],
\end{equation}
with dissipator
\begin{equation}
\mathcal{D}_{\mathrm{eff}}[\rho]
=
\sum_{n} \gamma_n
\left(
L_n \rho L_n^\dagger
- \frac{1}{2} \{ L_n^\dagger L_n, \rho \}
\right).
\end{equation}

The coefficients $\gamma_n$ scale as
\begin{equation}
\gamma_n \propto \frac{1}{M^4},
\end{equation}
consistent with the dimension-8 suppression.

This structure provides the microscopic origin of the thermodynamic arrow of
time: the visible sector behaves as an open quantum system after tracing out the
mirror sector, and the Lindblad dissipator generates entropy production.
% ========================================================
% 4 CONCLUSION
% ========================================================

\section{Conclusion}

This Companion develops the corrected thermodynamic sector of RDCC in its
v49.1 form. The visible and mirror sectors possess opposite thermodynamic
arrows of time, ensuring global CPT invariance. The mirror sector cannot
dissipate energy, leading to a thermodynamic blockade that prevents the
formation of compact objects and results in diffuse, non-compact halos.

The global entropy structure, the Lindblad-like effective dissipator, and the
dimension-8 inter-sector operator form the backbone of RDCC's thermodynamic
interpretation. These results integrate the relaxation sector (Companion III),
the ultraviolet sector (Companion X), the Boltzmann hierarchy (Companion VII),
and the growth suppression (Companion XIII).

\end{document}
